\begin{tikzpicture}[american, transform shape]

  % -------------------------------------------------------------
  % Entradas de señal (horizontales alineadas con los terminales +)
  % -------------------------------------------------------------
  % U1 con yscale=-1 tiene su terminal U1.+ en y = 2.99
  % U2 con yscale=1 tiene su terminal U2.+ en y = -2.99
  \node[ocirc] (in1) at (0.3, 2.99) {};
  \node[anchor=east] at (in1.west) {$e_1$};
  \node[ocirc] (in2) at (0.3, -2.99) {};
  \node[anchor=east] at (in2.west) {$e_2$};

  % -------------------------------------------------------------
  % Etapa 1: Buffers de entrada (U1 y U2)
  % -------------------------------------------------------------
  \node[op amp, yscale=-1] (U1) at (3.0, 2.5) {};
  \node at (2.8, 2.5) {$U_1$};
  \draw (in1) -- (U1.+);

  \node[op amp] (U2) at (3.0, -2.5) {};
  \node at (2.8, -2.5) {$U_2$};
  \draw (in2) -- (U2.+);

  % Nodos de terminales inversores
  \draw (U1.-) -- (1.8, 2.01);
  \draw (U2.-) -- (1.8, -2.01);

  % Impedancia Z1 entre terminales inversores
  \draw (1.8, 0.9) to[generic, l={$Z_1$}] (1.8, -0.9);
  \draw (1.8, 2.01) -- (1.8, 0.9);
  \draw (1.8, -2.01) -- (1.8, -0.9);

  % Impedancias de realimentacion Z2
  \draw (1.8, 1.3) node[circ]{} -- (2.3, 1.3)
        to[generic, l_={$Z_2$}] (3.7, 1.3)
        -- (4.2, 1.3) -- (4.2, 2.5) -- (U1.out);

  \draw (1.8, -1.3) node[circ]{} -- (2.3, -1.3)
        to[generic, l={$Z_2$}] (3.7, -1.3)
        -- (4.2, -1.3) -- (4.2, -2.5) -- (U2.out);

  % Tensiones intermedias eo1 y eo2
  \draw (U1.out) -- (4.8, 2.5) node[circ]{};
  \node[above] at (4.8, 2.5) {$e_{o1}$};
  \draw (U2.out) -- (4.8, -2.5) node[circ]{};
  \node[below] at (4.8, -2.5) {$e_{o2}$};

  % -------------------------------------------------------------
  % Etapa 2: Restador (U3)
  % -------------------------------------------------------------
  \node[op amp] (U3) at (9.0, 0.0) {};
  \node at (8.8, 0.0) {$U_3$};

  % Rama superior: eo1 -> Z3 -> entrada inversora U3.-
  \draw (4.8, 2.5) -- (5.0, 2.5)
        to[generic, l={$Z_3$}] (6.3, 2.5)
        -- (6.6, 2.5) |- (U3.-);

  % Rama inferior: eo2 -> Z3 -> entrada no inversora U3.+
  \draw (4.8, -2.5) -- (5.0, -2.5)
        to[generic, l_={$Z_3$}] (6.3, -2.5)
        -- (6.6, -2.5) |- (U3.+);

  % Realimentacion Z4 en lazo inversor (desde nodo antes de U3.- hacia U3.out)
  \draw (7.3, 0.49) node[circ]{} -- (7.3, 1.8)
        to[generic, l={$Z_4$}] (9.6, 1.8)
        -| (U3.out);
  \draw (9.6, 0.0) node[circ]{};

  % Impedancia Z4 a masa en lazo no inversor (desde nodo antes de U3.+ hacia masa)
  \draw (7.3, -0.49) node[circ]{}
        to[generic, l={$Z_4$}] (7.3, -1.9)
        node[ground]{};

  % Salida final eo
  \draw (U3.out) -- (10.6, 0.0);
  \node[ocirc] (eo) at (10.6, 0.0) {};
  \node[anchor=west] at (eo.east) {$e_o$};

\end{tikzpicture}
